\section{Method}
\label{sec:method}

We investigate how action-relevant representations are organized across three VLA architectures using complementary probing, attention, and perturbation analyses.

%------------------------------------------------------------------
\subsection{Architectures and Data}
\label{sec:arch_data}

\looseness=-1
We study three VLAs spanning distinct backbone families (Table~\ref{tab:arch_overview}):
\textbf{OpenVLA}~\cite{kim2024_openvla,karamcheti2024} (Llama-2-7B~\cite{touvron2023_llama2}, 32 layers, pretrained on Open X-Embodiment~\cite{padalkar2023_oxe} [OXE; 970K episodes]),
\textbf{TraceVLA}~\cite{zheng2024tracevla} (Phi-3-mini~\cite{abdin2024phi3}, 32 layers, trained on an OXE subset with spatial trace conditioning), and
\textbf{NORA}~\cite{hung2025nora} (Qwen2.5-VL-3B, 36 layers, grouped-query attention with 16 query heads).
Each is evaluated in trained and reference conditions.
A critical asymmetry distinguishes the references: OpenVLA's is an OXE-pretrained VLA (already producing actions), while TraceVLA and NORA references are pure VLMs (Phi-3-Vision, Qwen2.5-VL) with no action training; thus OpenVLA comparisons measure VLA-to-VLA transitions, while TraceVLA and NORA measure VLM-to-VLA transitions.
OpenVLA additionally provides Llama-2 base and random-weight conditions (\S\ref{sec:openvla_routing}).

We evaluate on LIBERO~\cite{liu2023_libero}, comprising four task suites (Goal, Spatial, Object, Long-Horizon) with 10 tasks and 50 demonstrations each.
LIBERO provides a controlled evaluation framework: all suites share the same simulated tabletop environment and robot embodiment, isolating instruction complexity as the primary experimental variable.
The four suites systematically vary instruction demands (Goal requires object identification, Spatial adds relative positioning, Object introduces distractors, and Long-Horizon chains multiple sub-tasks), enabling us to test whether routing patterns generalize across instruction complexity levels.
The three architectures and four suites form a $3 \times 4$ matrix; 10 of 12 cells are filled (NORA Object and Long-Horizon lack available checkpoints).

%------------------------------------------------------------------
\subsection{Representation Extraction}
\looseness=-1
From each architecture, we extract hidden states at every transformer layer from two token positions: (a)~\textbf{image\_mean}, the mean activation over image tokens, and (b)~\textbf{readout}, the hidden state at the final pre-action token, which can attend to all preceding context.
L0 denotes the embedding layer, L1--L$N$ the post-transformer hidden states ($N = 32$ for OpenVLA/TraceVLA, $36$ for NORA).
We cross training status (trained vs.\ reference) $\times$ instruction presence (with vs.\ without) across all 10 cells.
OpenVLA additionally provides DinoSigLIP projector, DINOv2, and CLIP representations for the single-architecture analyses (\S\ref{sec:openvla_routing}).

\begin{figure}[t]
\centering
\includegraphics[width=0.8\columnwidth]{figures/fig_rsa_perlayer_curve.pdf}
\caption{Per-layer RSA (Spearman $\rho$) for trained/reference OpenVLA. Trained: inverted-U peaking at L12; reference: monotonic decline. Readout and image\_mean diverge after L12.}
\label{fig:rsa_perlayer}
\end{figure}

%------------------------------------------------------------------
\subsection{Cross-Architecture Analyses}
\label{sec:cross_arch_analyses}

Three analyses are applied to every available cell in the $3 \times 4$ matrix.

\paragraph{Action probes.}\looseness=-1
Two-layer MLP probes~\cite{alain2017} predict 7 degree-of-freedom (DoF) actions from frozen representations (details in Appendix~\ref{sec:openvla_methods_appendix}).
Probe $R^2$ measures decodability, not causal usage.
$\Delta R^2 = R^2_{\text{peak}} - R^2_{\text{exit}}$; we use $> 0.10$ as the threshold for substantial peak-to-exit reduction.
Demo-level train/test splits and Monte Carlo cross-validation ($n = 1{,}000$) provide 95\% confidence intervals (CIs; Appendix~\ref{sec:delta_r2_threshold}); linear probes yield consistent patterns (Appendix~\ref{sec:linear_probe}).

\paragraph{L1 attention head profiling.}\looseness=-1
To determine whether attention heads preferentially route instruction information to the readout position, each L1 head is classified as instruction-dominant ($>50\%$ of readout-token attention to instruction tokens) or image-dominant, reported as $k/N$ heads (32 for MHA, 16 for GQA).
Sensitivity analysis at thresholds from $30\%$ to $70\%$ confirms ordering robustness (Appendix~\ref{sec:threshold_sensitivity}).

\paragraph{Counterfactual instruction perturbation.}\looseness=-1
To test whether the model's actions depend on instruction content, the instruction is replaced with four conditions: empty string, wrong task, shuffled tokens, and paraphrased.
We measure trained-model sensitivity via Cohen's $d$ on action MSE (OpenVLA, TraceVLA) or cosine distance (NORA); reference VLMs use cosine distance on hidden representations.
These metrics operate on different scales; cross-metric ratios are invalid (\S\ref{sec:cf_sensitivity}).

\looseness=-1
OpenVLA additionally supports cross-instruction probes, representational similarity analysis (RSA)~\cite{kriegeskorte2008}, and external baselines (DINOv2~\cite{oquab2024_dinov2}, CLIP~\cite{radford2021_clip}); these single-architecture methods are detailed in Appendix~\ref{sec:openvla_methods_appendix}.
Figure~\ref{fig:rsa_perlayer} previews this: per-layer geometric alignment ($\rho$) diverges between trained and reference models after L12, with trained readout tokens maintaining alignment ($\rho = 0.584$) while image-mean representations decay ($\rho = 0.422$), suggesting task-specific conditioning at the readout position.

\begin{figure}[t]
\centering
\includegraphics[width=\columnwidth]{figures/fig_matrix_heatmap.pdf}
\caption{\textbf{Cross-architecture analysis matrix.} Three architectures (OpenVLA, TraceVLA, NORA) $\times$ four LIBERO suites; 10 of 12 cells filled. Color encodes $\Delta R^2$ (peak$-$exit); bold borders mark VLM-to-VLA transitions exhibiting substantial exit collapse ($\Delta R^2 > 0.10$).}
\label{fig:result_matrix}
\end{figure}

\FloatBarrier
